\chapter{Formulaire}

\noindent\textit{Toutes les formules du programme de 3ème, regroupées par thème pour une
révision rapide. Chaque formule est démontrée et illustrée dans le chapitre correspondant :
ce formulaire n'est qu'un aide-mémoire, pas un cours.}
\vspace{8pt}

\begin{multicols}{2}

\subsection*{\color{onbo}Puissances}
\begin{itemize}
\item $a^{m}\times a^{n}=a^{m+n}$ \quad $\dfrac{a^{m}}{a^{n}}=a^{m-n}$
\item $(a^{m})^{n}=a^{m\times n}$ \quad $(a\times b)^{n}=a^{n}\times b^{n}$ \quad
$\left(\dfrac ab\right)^{n}=\dfrac{a^{n}}{b^{n}}$
\item $a^{-n}=\dfrac1{a^{n}}$ \quad $a^{0}=1$ (pour $a\ne0$)
\item Notation scientifique : $a\times10^{n}$ avec $1\le|a|<10$
\end{itemize}

\subsection*{\color{onbo}Racines carrées}
\begin{itemize}
\item $\left(\sqrt a\right)^{2}=a$ et $\sqrt{a^{2}}=a$ (pour $a\ge0$)
\item $\sqrt{a\times b}=\sqrt a\times\sqrt b$ \quad
$\sqrt{\dfrac ab}=\dfrac{\sqrt a}{\sqrt b}$ (pour $a\ge0,b>0$)
\item $a<b\iff\sqrt a<\sqrt b$ (pour $a,b\ge0$)
\item Rendre rationnel : $\dfrac{c}{\sqrt b}=\dfrac{c\sqrt b}{b}$ (pour $b>0$)
\item $x^{2}=a$ : deux solutions $\pm\sqrt a$ si $a>0$ ; une solution $0$ si $a=0$ ;
aucune solution si $a<0$
\end{itemize}

\subsection*{\color{onbo}Identités remarquables}
\begin{itemize}
\item $(a+b)^{2}=a^{2}+2ab+b^{2}$
\item $(a-b)^{2}=a^{2}-2ab+b^{2}$
\item $(a-b)(a+b)=a^{2}-b^{2}$
\item Distributivité : $k(a+b)=ka+kb$ \quad $(a+b)(c+d)=ac+ad+bc+bd$
\end{itemize}

\subsection*{\color{onbo}Arithmétique}
\begin{itemize}
\item PGCD par l'algorithme d'Euclide : divisions en cascade, le dernier reste non nul
\item $\text{PGCD}(a\,;b)\times\text{PPCM}(a\,;b)=a\times b$
\item Fraction irréductible : diviser $a$ et $b$ par $\text{PGCD}(a\,;b)$
\end{itemize}

\subsection*{\color{onbo}Proportionnalité \& pourcentages}
\begin{itemize}
\item Produit en croix : $\dfrac ab=\dfrac cd\iff a\times d=b\times c$
\item Taux de variation : $t=\dfrac{V_f-V_i}{V_i}\times100$ (en $\%$)
\item Augmenter de $t\,\%$ : $\times\left(1+\dfrac t{100}\right)$ ; diminuer de $t\,\%$ :
$\times\left(1-\dfrac t{100}\right)$
\item Deux évolutions successives : on \textbf{multiplie} les coefficients
\end{itemize}

\subsection*{\color{onbo}Équations \& systèmes}
\begin{itemize}
\item Produit nul : $A\times B=0\iff A=0$ ou $B=0$
\item Inéquation : multiplier/diviser par un nombre négatif \textbf{inverse le sens}
\item Système $ax+by=c$ : substitution, combinaison, ou intersection de droites
\end{itemize}

\subsection*{\color{onbo}Fonctions affines}
\begin{itemize}
\item Fonction affine : $f(x)=ax+b$ ($a$ = coefficient directeur, $b$ = ordonnée à
l'origine)
\item Par deux points $A(x_1\,;y_1)$, $B(x_2\,;y_2)$ : $a=\dfrac{y_2-y_1}{x_2-x_1}$,
$b=y_1-a\,x_1$
\item Par un point $(x_0\,;y_0)$ et $a$ connu : $b=y_0-a\,x_0$
\end{itemize}

\subsection*{\color{onbo}Théorème de Thalès}
\begin{itemize}
\item $D\in[AB]$, $E\in[AC]$, $(DE)\parallel(BC)$ : $\dfrac{AD}{AB}=\dfrac{AE}{AC}=\dfrac{DE}{BC}$
\item Réciproque : $\dfrac{AD}{AB}=\dfrac{AE}{AC}\Rightarrow(DE)\parallel(BC)$
\item Coefficient $k$ : longueurs $\times k$, périmètre $\times k$, aire $\times k^{2}$
\end{itemize}

\subsection*{\color{onbo}Théorème de Pythagore}
\begin{itemize}
\item Triangle rectangle en $A$ : $BC^{2}=AB^{2}+AC^{2}$ (hypoténuse $BC$)
\item Réciproque : $BC^{2}=AB^{2}+AC^{2}\Rightarrow$ rectangle en $A$
\item Distance dans un repère : $AB=\sqrt{(x_B-x_A)^{2}+(y_B-y_A)^{2}}$
\end{itemize}

\subsection*{\color{onbo}Trigonométrie}
\begin{itemize}
\item $\cos\widehat B=\dfrac{\text{adjacent}}{\text{hypoténuse}}$ \quad
$\sin\widehat B=\dfrac{\text{opposé}}{\text{hypoténuse}}$ \quad
$\tan\widehat B=\dfrac{\text{opposé}}{\text{adjacent}}$
\item $\cos^{2}\widehat B+\sin^{2}\widehat B=1$ \quad
$\tan\widehat B=\dfrac{\sin\widehat B}{\cos\widehat B}$
\end{itemize}

\subsection*{\color{onbo}Vecteurs}
\begin{itemize}
\item Relation de Chasles : $\vec{AB}+\vec{BC}=\vec{AC}$
\item Coordonnées : $\vec{AB}\binom{x_B-x_A}{y_B-y_A}$ \quad
$\|\vec{AB}\|=\sqrt{(x_B-x_A)^{2}+(y_B-y_A)^{2}}$
\item $ABCD$ parallélogramme $\iff\vec{AB}=\vec{DC}$
\end{itemize}

\subsection*{\color{onbo}Angle inscrit, angle au centre}
\begin{itemize}
\item $\widehat{AOB}=2\times\widehat{ACB}$ (angle au centre = double de l'angle inscrit)
\item Angles inscrits interceptant le même arc : \textbf{égaux}
\item $[AB]$ diamètre, $C$ sur le cercle $\Rightarrow ABC$ rectangle en $C$ (et réciproque)
\end{itemize}

\subsection*{\color{onbo}Aires et périmètres (rappels)}
\begin{itemize}
\item Carré de côté $c$ : $P=4c$, $\mathcal A=c^{2}$
\item Rectangle $L\times\ell$ : $P=2(L+\ell)$, $\mathcal A=L\times\ell$
\item Triangle, base $b$, hauteur $h$ : $\mathcal A=\dfrac{b\times h}2$
\item Disque de rayon $r$ : $P=2\pi r$, $\mathcal A=\pi r^{2}$
\end{itemize}

\subsection*{\color{onbo}Solides usuels}
\begin{itemize}
\item Prisme droit / cylindre : $V=\mathcal A_{\text{base}}\times h$ ;
$\mathcal A_{\text{lat}}=P_{\text{base}}\times h$ (cylindre : $2\pi rh$, $V=\pi r^{2}h$)
\item Pyramide : $V=\dfrac13\mathcal A_{\text{base}}\times h$ ; si régulière,
$\mathcal A_{\text{lat}}=\dfrac12P_{\text{base}}\times a$ ($a$ = apothème)
\item Cône : $l^{2}=r^{2}+h^{2}$ ; $\mathcal A_{\text{lat}}=\pi rl$ ;
$\mathcal A_{\text{tot}}=\pi r(r+l)$ ; $V=\dfrac13\pi r^{2}h$
\item Sphère / boule : $\mathcal A=4\pi r^{2}$ ; $V=\dfrac43\pi r^{3}$
\item Section parallèle à la base, coefficient $k$ : longueurs $\times k$, aires
$\times k^{2}$, volumes $\times k^{3}$
\end{itemize}

\subsection*{\color{onbo}Statistiques}
\begin{itemize}
\item Fréquence $=$ effectif $\div$ effectif total
\item Moyenne : $\bar x=\dfrac{\sum(\text{valeur}\times\text{poids})}{\sum\text{poids}}$
\item Médiane (série ordonnée) : $N$ impair $\Rightarrow$ position $\dfrac{N+1}2$ ; $N$
pair $\Rightarrow$ moyenne des positions $\dfrac N2$ et $\dfrac N2+1$
\item Étendue $=$ plus grande valeur $-$ plus petite valeur
\item Diagramme circulaire : angle $=$ fréquence $\times360^{\circ}$
\end{itemize}

\end{multicols}
